% De la tela a la rejilla — los cuatro pilares del pensamiento computacional
% vistos en el calado cartagüeño.
%
% Diagrama propio en TikZ (sin librerías extra: las tramas se dibujan con
% \foreach y líneas simples). Se tiñe con el color de la línea
% (milcGradeDark = verde lima de Pensamiento computacional).
%
% Muestra el mecanismo real que describe Pérez-Bustos (2019): la tela llega
% densa; la caladora marca un cuadro y saca los hilos DE UNO EN UNO hasta
% dejar una rejilla regular; sobre esa rejilla —y no sobre la tela— se teje
% la puntada que se repite. Debajo, el orden fijo de los pasos.
%
% Las anotaciones de dentro enseñan (nombran el pilar en cada tránsito); la
% síntesis va en el pie de figura vía \guiaDiagrama, no aquí.
\begin{tikzpicture}[scale=1,font=\sffamily,>=stealth]

  % ── Panel 1 · la tela densa: todos los hilos, nada distinguible ──────
  \begin{scope}[shift={(0,0)}]
    \draw[milcNegro,line width=0.6pt] (0,0) rectangle (2.6,2.6);
    \foreach \i in {0.1,0.2,...,2.5}{
      \draw[milcNegro!45,line width=0.25pt] (\i,0) -- (\i,2.6);
      \draw[milcNegro!45,line width=0.25pt] (0,\i) -- (2.6,\i);
    }
    \node[below,font=\scriptsize\bfseries,milcNegro] at (1.3,-0.2) {la tela};
    \node[below,font=\scriptsize\itshape,milcNegro!75,text width=2.8cm,align=center]
      at (1.3,-0.5) {urdimbre y trama: todo junto, nada legible};
  \end{scope}

  % ── Panel 2 · la rejilla: hilos retirados de uno en uno ──────────────
  \begin{scope}[shift={(4.2,0)}]
    \draw[milcNegro,line width=0.6pt] (0,0) rectangle (2.6,2.6);
    \foreach \i in {0.1,0.2,...,2.5}{
      \draw[milcNegro!25,line width=0.2pt] (\i,0) -- (\i,2.6);
      \draw[milcNegro!25,line width=0.2pt] (0,\i) -- (2.6,\i);
    }
    % los hilos que quedan tras el deshilado (rejilla regular 4x4)
    \foreach \i in {0.3,0.9,1.5,2.1}{
      \draw[milcGradeDark,line width=1.1pt] (\i,0.2) -- (\i,2.4);
      \draw[milcGradeDark,line width=1.1pt] (0.2,\i) -- (2.4,\i);
    }
    \node[below,font=\scriptsize\bfseries,milcNegro] at (1.3,-0.2) {la rejilla};
    \node[below,font=\scriptsize\itshape,milcNegro!75,text width=2.8cm,align=center]
      at (1.3,-0.5) {un hilo a la vez, contados: la trama se vuelve visible};
  \end{scope}

  % ── Panel 3 · la puntada que se repite sobre la rejilla ──────────────
  \begin{scope}[shift={(8.4,0)}]
    \draw[milcNegro,line width=0.6pt] (0,0) rectangle (2.6,2.6);
    \foreach \i in {0.3,0.9,1.5,2.1}{
      \draw[milcNegro!30,line width=0.7pt] (\i,0.2) -- (\i,2.4);
      \draw[milcNegro!30,line width=0.7pt] (0.2,\i) -- (2.4,\i);
    }
    % la misma puntada, repetida en cada celda de la rejilla
    \foreach \x in {0.6,1.2,1.8}{
      \foreach \y in {0.6,1.2,1.8}{
        \draw[milcGradeDark,line width=0.9pt]
          (\x-0.22,\y-0.22) -- (\x+0.22,\y+0.22)
          (\x-0.22,\y+0.22) -- (\x+0.22,\y-0.22);
        \draw[milcMostaza,line width=0.9pt] (\x,\y) circle (0.14);
      }
    }
    \node[below,font=\scriptsize\bfseries,milcNegro] at (1.3,-0.2) {la puntada};
    \node[below,font=\scriptsize\itshape,milcNegro!75,text width=2.8cm,align=center]
      at (1.3,-0.5) {una sola forma, repetida: el diseño aparece};
  \end{scope}

  % ── Tránsitos: cada flecha nombra el pilar que opera ─────────────────
  \draw[->,milcGradeDark,line width=1.2pt] (2.9,1.3) -- (3.9,1.3);
  \node[above,font=\scriptsize\bfseries,milcGradeDark,text width=1.4cm,align=center]
    at (3.4,1.45) {descom-\\poner};

  \draw[->,milcGradeDark,line width=1.2pt] (7.1,1.3) -- (8.1,1.3);
  \node[above,font=\scriptsize\bfseries,milcGradeDark,text width=1.5cm,align=center]
    at (7.6,1.45) {abstraer};

  \node[font=\scriptsize\bfseries,milcMostaza!80!black] at (5.5,3.0) {patrón};
  \draw[milcMostaza,line width=0.9pt] (4.4,2.85) -- (5.0,2.85);
  \draw[milcMostaza,line width=0.9pt] (6.0,2.85) -- (6.6,2.85);

  % ── El orden fijo: cambiarlo arruina la pieza ────────────────────────
  \node[font=\scriptsize\bfseries,milcGradeDark] at (0.9,-1.95) {algoritmo};
  \foreach \i/\x/\texto in {%
    1/0.9/{marcar el cuadro},%
    2/3.6/{sacar 2 y 2},%
    3/6.3/{deshilar de uno en uno},%
    4/9.0/{remendar y bordar}}{
    \node[fill=milcGradeDark,text=white,rounded corners=3pt,inner sep=3pt,
          font=\scriptsize\bfseries,text width=2.0cm,align=center] (P\i) at (\x,-2.5) {\i. \texto};
  }
  \foreach \a/\b in {1/2,2/3,3/4}{
    \draw[->,milcGradeDark,line width=1pt] (P\a) -- (P\b);
  }

\end{tikzpicture}
